% Part: first-order-logic
% Chapter: beyond
% Section: other-logics

\documentclass[../../../include/open-logic-section]{subfiles}

\begin{document}

\olfileid{fol}{byd}{oth}

\olsection{منطق‌های دیگر}

احتمالاً تا اینجا دریافته‌اید که طراحی منطقی تازه کار دشواری نیست.
شما نیز می‌توانید نحو خود را بسازید، دستگاهی استنتاجی ابداع کنید و
معناشناسی متناسبی برایش بپردازید. اگر بخواهید دستگاه !!{derivation}تان
نسبت به معناشناسی تام باشد، شاید لازم باشد قدری هوشمندی به خرج دهید؛ و
شاید قانع‌کردن عموم مردم به اینکه منطق شما حقیقتاً جالب است اندکی تلاش
بخواهد. اما در عوض می‌توانید ساعت‌ها سرگرمی سالم و دلپذیر داشته باشید
و خواص ریاضی و محاسباتی منطقتان را کاوش کنید.

دهه‌های اخیر شاهد انفجاری واقعی در منطق‌های صوری بوده‌اند. منطق فازی
برای بازنمایی استدلال دربارهٔ خواص مبهم طراحی شده است. منطق احتمالاتی
برای بازنمایی استدلال دربارهٔ عدم قطعیت طراحی شده است. منطق‌های پیش‌فرض
و منطق‌های نامونوتون برای بازنمایی صورت‌های ابطال‌پذیر استدلال طراحی
شده‌اند؛ یعنی استنتاج‌های «معقولی» که ممکن است بعداً با پدیدارشدن
اطلاعات تازه کنار گذاشته شوند. منطق‌های معرفتی برای بازنمایی استدلال
دربارهٔ معرفت، منطق‌های علّی برای بازنمایی استدلال دربارهٔ روابط علّی،
و حتی منطق‌های «تکلیفی» برای بازنمایی استدلال دربارهٔ الزام‌های اخلاقی
طراحی شده‌اند. بسته به اینکه انگیزهٔ اصلی معرفی این دستگاه‌ها فلسفی،
ریاضی یا محاسباتی باشد، ممکن است چنین موجوداتی را زیر عنوان منطق
ریاضی، منطق فلسفی، هوش مصنوعی، علوم شناختی یا در جایی دیگر مطالعه‌شده
بیابید.

این فهرست همچنان ادامه دارد و امکان‌ها پایان‌ناپذیر می‌نمایند. شاید
هرگز به رؤیای لایبنیتس، یعنی فروکاستن همهٔ عقل انسانی به محاسبه، دست
نیابیم---اما این نمی‌تواند ما را از کوشش بازدارد.

\end{document}
